\chapter{盛世帝国：雍正、乾隆与边疆经营}
\label{ch:high-qing}
\begin{figure}[htbp]
  \centering
  \begin{tikzpicture}[
    x=1cm,y=1cm,
    date/.style={font=\scriptsize\bfseries, text=dynasty},
    event/.style={draw=timeline, fill=paper, rounded corners=2pt, align=center,
      inner sep=3pt, font=\scriptsize, text=ink},
    note/.style={font=\scriptsize\color{source}, align=center}
  ]
    \draw[timeline, line width=1.2pt] (0,0) -- (12,0);
    \foreach \x/\year in {0/1723,1.6/1728,3.0/1735,5.0/1755,6.7/1759,8.3/1768,10.1/1786,12/1795}
      {\draw[timeline] (\x,.12) -- (\x,-.12); \node[date, below] at (\x,-.18) {\year};}
    \node[event, above] at (0,.42) {雍正即位；\\财政与西北军务};
    \node[event, below] at (1.6,-.62) {准噶尔远征、\\驻藏部署};
    \node[event, above] at (3.0,.42) {乾隆即位；\\山地与边地压力};
    \node[event, below] at (5.0,-.62) {伊犁战事\\开始};
    \node[event, above] at (6.7,.42) {天山南路\\战事收束};
    \node[event, below] at (8.3,-.62) {第二次金川\\战争升级};
    \node[event, above] at (10.1,.42) {台湾起事与\\廓尔喀边境战事};
    \node[event, below] at (12,-.62) {禅位；嘉庆\\承接财政与战事压力};
    \node[note] at (6,-1.35)
      {军行、设制、迁调与地方治理是不同阶段；这条时间线不把它们压缩为“扩张完成”。};
  \end{tikzpicture}
  \caption{雍正、乾隆时期的编年定位（1723--1795，简表）。图以财政--军政、伊犁与天山南路、金川、台湾和高原的关键节点标示本章阅读顺序。主要依《清世宗实录》《清高宗实录》、军机处档和相关方略；方略的军功、兵数与“平定”语言须与地方、多语和物质材料互证。}
  \label{fig:high-qing-timeline}
\end{figure}

\begin{figure}[htbp]
  \centering
  \begin{tikzpicture}[
    x=1cm,y=1cm,
    place/.style={draw=timeline, fill=paper, rounded corners=2pt, align=center,
      inner sep=3pt, font=\small},
    court/.style={place, draw=dynasty, fill=dynasty!13},
    frontier/.style={place, draw=timeline, fill=timeline!13},
    local/.style={place, draw=source, fill=source!10},
    note/.style={align=center, font=\scriptsize\color{source}}
  ]
    \fill[paper] (-.7,-3.85) rectangle (7.8,3.85);
    \draw[dynasty, line width=.8pt] (-.7,3.85) -- (7.8,3.85);
    \node[font=\bfseries\large, text=dynasty] at (3.55,3.42)
      {1755--1765：伊犁、绿洲与高原的军政节点};

    \node[court] (beijing) at (3.65,2.55) {北京\\军机、理藩院与\\军需文书};
    \node[frontier] (mongolia) at (1.15,1.55) {蒙古诸部\\盟旗与军行};
    \node[frontier] (ili) at (-.05,.1) {伊犁\\1755后战事与\\1760将军设置};
    \node[frontier] (urumqi) at (2.05,-.15) {迪化、北疆\\驻防与屯田};
    \node[local] (oases) at (4.45,-.2) {喀什噶尔、叶尔羌\\1758--1759战事\\与绿洲社会};
    \node[frontier] (qinghai) at (.85,-2.1) {青海、川西\\道路、牲畜与\\转运压力};
    \node[local] (lhasa) at (3.1,-2.8) {拉萨、寺院网络\\驻藏制度调整};
    \node[local] (nepal) at (5.75,-2.55) {喜马拉雅方向\\1780年代后战事};
    \node[local] (jinchuan) at (6.6,1.45) {金川山地\\征调、道路与\\地方中介};

    \draw[dynasty, dashed] (beijing) -- node[above,sloped,note]
      {文书与调度（示意）} (mongolia);
    \draw[->, dynasty, line width=1.1pt] (mongolia) -- node[above,sloped,note]
      {分路军行} (ili);
    \draw[->, timeline, line width=1.1pt] (ili) -- node[above,note]
      {驻防、迁调与贸易节点} (urumqi);
    \draw[->, timeline, line width=1.1pt] (urumqi) -- node[above,note]
      {天山南北路的不同条件} (oases);
    \draw[->, dynasty, dashed] (beijing) to[out=-145,in=80] node[left,note]
      {高原军政联系（示意）} (lhasa);
    \draw[->, timeline, dashed] (qinghai) -- node[below,sloped,note]
      {道路、寺院与地方协商} (lhasa);
    \draw[timeline, dashed] (lhasa) -- node[below,sloped,note]
      {边防、使节与后续冲突} (nepal);
    \draw[source, dashed] (beijing) to[out=-25,in=75] node[right,note]
      {山地军需并非直线可达} (jinchuan);

    \node[draw=source, fill=paper, rounded corners=3pt, align=center,
      text=source, font=\small] at (5.8,2.45)
      {节点与箭头只提示可见的\\军行、文书和交通关系；\\不画连续疆界或稳定控制面};
  \end{tikzpicture}
  \caption{高宗前中期的西北与高原军政关系（1755--1765为主，示意）。图以伊犁战事、北疆驻防、南疆绿洲、青海--拉萨交通和金川军需等不同性质的节点组织，提示征服、驻防与日常治理并非同一过程。依据《清高宗实录》乾隆二十年至三十年条、《平定准噶尔方略》《平定回部方略》、军机处军需与将领奏折、理藩院满汉蒙藏档；绿洲部分须与维吾尔／察合台文文书、地方材料及俄国边境档并读，西藏部分须与藏文寺院和地方材料互证。路线与位置均为约略示意，不按比例；箭头不表示对牧地、绿洲、寺院、山地或边界的持续主权和社会接受，方略所载军数、伤亡与“平定”亦不能视作无须核验的总数。}
  \label{fig:high-qing-frontier-ili-tibet}
\end{figure}


\section{1723—1735：把亏空、密折与西北放到同一张案上}

雍正元年，继承的并不是一座已经安静的北京城，而是一串仍待结算的账目、宗室关系和西北军情。\eventanchor{EQ-1723-YONGZHENG-ACCESSION}{雍正即位（1723）}后，皇帝以朱批和用人把来自各省、内廷与边地的消息拉近决策中心；年羹尧、隆科多先后失势，则显示这种收拢同时在重划近臣和将领可以占据的位置。宫中档案能让人看见诏令、奏折和处分如何往返，却不能把朝廷的猜忌直接写成所有官员共有的动机。

雍正四年推进的\eventanchor{EQ-1726-HAOXIAN-GUIGONG}{耗羡归公与养廉银（1726）}，把原先散在地方征收中的附加收入试图纳入较可核查的财政秩序。它既关乎州县官能否有合法的办事经费，也关乎户部能否看清亏空；但一纸成例并不等于每县的征收、折色或摊派随即相同。三年后设置军机房，正是在准噶尔战事催促下为紧急文书另开较短的通道。文书走得更快，驼马、粮草与道路却未必同样迅速抵达草原。1731、1732年西北远征的受挫，把这个差距暴露出来：军机房解决的是消息与批答的路径，不是高原和荒漠中补给的距离。

\begin{turningpoint}[军机房：缩短批答路径，不能缩短草原距离（1729）]
\eventanchor{EQ-1729-JUNJICHU-FOUNDING}{军机房的设置（1729）}改变了紧急军政文书进入皇帝决策的路径；它并未替西北军队解决驼马、粮草、水源和道路。制度的速度与补给的速度必须分开衡量。
\end{turningpoint}

这条分别也贯穿高原事务。雍正六年在准噶尔威胁下\eventref{EQ-1728-ZHUNGHAR-TIBET-SETTLEMENT}{调整驻藏与青海蒙古的军事部署（1728）}，使拉萨、青海和北京之间的军政联络更紧密；但命令仍须经过盟旗关系、寺院网络、驿道和运输者。把军机房、养廉银和驻藏部署并置，并不是要把它们说成同一项政策，而是要看到一项跨区域治理的共同难题：中枢要使钱、信和人到达不同地方，所倚赖的总是当地能够协商、转运和供给的关系。账册能记录报解与核销，不能替我们丈量一段山道的可通行性，更不能把地方中介写成被动的传声筒。

\section{1735—1754：继承一套制度，也继承它的边缘}

1735年乾隆即位，接手的是已被整顿的财政和军机决策机制，也是贵州苗疆、川西山地、青海与准噶尔边缘持续传来的压力。荣江一带的冲突、驻兵和设治，不能被“平定”的结论吞没；对于村寨、迁徙和暴力，地方档案与后出的地方记忆往往比实录更接近问题，却也不能据此推定整个苗疆的经历。1738年至1746年的第一次金川战争同样证明，山地道路、土司联盟、粮运和招抚并非战报的背景。和议使战事暂告收束，却没有消除中央在川西经营所需的中介和成本。

河工、盐政、矿务和赈济在这些年持续占用财政注意力。乾隆朝初年的“宽大”诏令、旗籍调整和灾赈措置，分别处理政治案件、供养资格和粮价压力；它们不应被压缩为一条单向的集权线。1730年代以来的奏销、朱批和户部档可以追到国家怎样界定问题，若要知道某一河段的劳役、某一矿区的流动人口或某一旗人家庭的债务，仍须回到地方账册、契约和个案材料。

宗教与知识的处理也不沿着一条从禁到用、或从开放到封闭的直线前进。雍正二年\eventref{EQ-1724-CATHOLIC-PROSCRIPTION}{扩大对天主教活动的限制（1724）}，把传教、地方教民与外来人员一并纳入地方治理的视野；到1744年，宫廷仍\eventanchor{EQ-1744-JESUIT-CALENDAR-SERVICE}{使用部分西洋传教士的历算与技术服务（1744）}。前者是限制的命令，后者是宫廷知识需求留下的记录，不能互相抵销，也不能推出所有地区的宗教生活。地方官在户口、治安、礼仪和外来交涉之间如何处置，教民与传教士又如何调整活动，须在各地案件、书信和教会材料中逐处追问。国家可以同时限制某些宗教网络、征用其中一部分知识；这正使“国家意志”不能代替具体关系的历史。

\section{1755—1777：伊犁、金川与一座扩大的帝国}

1754年阿睦尔撒纳归附改变了清准竞争的联盟格局。次年，\eventanchor{EQ-1755-ILI-EXPEDITION}{清军分路进入伊犁（1755）}，击溃达瓦齐；阿睦尔撒纳旋即反清，1758—1759年又须越过天山南路处理大小和卓的政权。北京收到的方略把分路进军、奖惩和凯旋编为清楚的次序；草原和绿洲中的盟约、疾病、逃散与军粮，不能只从这部胜利叙事中复原。1760年设伊犁将军，驻防、屯田和安置随之展开，意味着征服转化为长期的水源、牧地、军饷和户口安排，而非一条疆界线的落定。

乌什1765年的反清事件提醒人们，驻防城和行政名目并不自动生成服从。西南的情形也相似：第二次金川战争自1768年延至1776年，炮兵、道路、民夫和四川及邻省的转运被拖入一场格外昂贵的山地战争。战后改土归流、设治和驻防改变了官府的接口，却没有替居民结清迁徙、土地和债务的代价。帝国在伊犁与金川获得了更大的调度范围，也把更多地方社会卷进它的供给链。

1757年，\eventanchor{EQ-1757-CANTON-TRADE-RESTRICTION}{西洋贸易被集中于广州口岸（1757）}，是同一时代另一种把流动纳入可管理地点的尝试。沿海治安、关税和交涉的顾虑，使官员、行商、船只与外来商人被放进港口的制度框架；它没有让海岸社会停止移动，也没有把贸易所得、走私风险或地方接触化为一笔全国通用的账。伊犁的军台、金川的山道和广州的码头面对的资源不同，却都提醒人们：制度只划定了应由谁呈报、征收或接待，真正的运输、信用和协商仍发生在具体的道路与水面上。

战争越向地方深处延伸，书写也越成为治理的一部分。1768年的\eventanchor{EQ-1768-LITERARY-INQUISITION}{文字狱案件（1768）}让著述、地方士人和政治语言进入案件链；1772年下令\eventanchor{EQ-1772-SIKU-QUANSHU-ORDER}{征集、编纂《四库全书》（1772）}，至1775年，\eventanchor{EQ-1775-SIKU-COLLECTION}{征书与审查在各省展开（1775）}。征集可以保存、整理和调动书籍，审查也可能使某些著作与作者面临风险，两者在同一网络中发生。禁毁目录、题本和编纂记录首先显示朝廷怎样分类文本，并不能直接告诉我们一本书在何处仍被抄写、转手或阅读；地方藏书、书信和幸存抄本也只能照见局部的流通。知识制度的力量由此可见，却不能被误写成对读者心意的完全控制。

\section{1778—1795：扩张的余波与治理的成本}

1770年代以后，土尔扈特部回归后的安置、《四库全书》的征集与审查、台湾林爽文事件、廓尔喀方向的战事和广州贸易，发生在不同的道路与制度场域。它们共同要求国家识别人口、调拨粮饷、处置边地关系，或规定港口中商人和外来船只的活动。1780年代台湾的战事尤其说明，海峡、垦殖、租佃和地方武装使“恢复控制”成为一个漫长的行政过程；1793年马戛尔尼使团的礼仪交涉也不能被倒写成后来战争的必然预演。

1786年\eventanchor{EQ-1786-LIN-SHUANGWEN-UPRISING}{林爽文起事在台湾爆发（1786）}，次年福康安等率军渡海，\eventanchor{EQ-1787-TAIWAN-CAMPAIGN}{海峡运输成为镇压行动的关键条件（1787）}。军报可以追出调兵、筹饷与收复城镇的次序，却不能以此抹平垦殖者、佃户、地方武装、港口商人和内山社群所处的不同位置。岛上的冲突并非由某一项赋役或某一种身份单独造成，镇压也不等于一切地方关系已经复原；府县档案、契约、航运记录和保存稀少的原住民材料所能提供的，仍是彼此不能替代的局部视野。

廓尔喀战争则把这一问题带回高原。1790年，清廷\eventanchor{EQ-1790-IMPERIAL-REGULATIONS-TIBET}{提出整顿西藏事务的制度性安排（1790）}；两年后，\eventanchor{EQ-1792-GURKHA-SETTLEMENT}{战后使节、边防与西藏军政的善后（1792）}成为新的问题。条文能够标出驻藏大臣和金瓶掣签等机制被强化的方向，不能预先规定寺院、噶厦、商路与边境来往会怎样回应。藏文寺院文书、清廷档案和跨喜马拉雅材料各自带着不同的机构语言；把其中一种叙述升格为整个高原的现实，会重新制造这些材料本来应当校正的盲点。

同一年，帝国还须在北京面对一条由广州延伸而来的海路。马戛尔尼使团在1793年的\eventanchor{EQ-1793-MACARTNEY-MISSION}{贸易、礼仪与外交交涉（1793）}，不是一场足以预告后来冲突的文化考试，而是英国扩张性的贸易要求与清廷既有港口、礼仪和接待秩序的一次相遇。中英双方留下的陈述都服务交涉本身；礼仪争执既不能证明一个完全不动的“闭关”帝国，也不能删去广州商人、官员和航运网络长期维持的实际往来。

乾隆末年留下的不是一个静止的“盛世”，而是一套可以调动档案、军队和多地资源、同时日益依赖地方执行与财政承受力的帝国。下一章从1796年的白莲教战争进入，观察这些负担怎样在战区、漕运、旗人生计和沿海贸易中逐步显形。

\section{档案能看见的路，档案看不见的人}

《清世宗实录》《清高宗实录》、朱批奏折、户部奏销、军机处与理藩院档案，使我们能够较清楚地追随报告抵达北京、命令离开北京的路线；《平定准噶尔方略》一类汇编则把军功与奖惩组织成朝廷能够认可的战争次序。它们最有力地证明的是文书、判断和程序曾经发生，而不是命令已在每一处落实。军费、伤亡、迁徙和安置尤其不能从一份凯旋奏报或一纸定额推成精确的社会总数。

因此，本章把伊犁、西藏、青海、川西、苗疆、台湾与广州放在同一条年代线上，却不把它们做成同一类地方。满、蒙、藏、维吾尔等材料、地方文书、契约、寺院记录以及俄国和邻国档案，能让边地、港口和地方社会不只作为北京文书的终点出现；它们的存留、译读和地点也不均衡，不能补成一张完整的帝国全景。高峰时期的可见力量正在这种不完整中显出边界：国家能更快地记账、发令和调兵，却始终要通过不同语言、道路、家庭与社群，才可能成为地方的经验。

\section{把钱粮、命令与远征扣在一起}

雍正元年，继承已经完成，康熙晚年留下的亏空、宗室供养、吏治和边务却仍要逐件结算。新朝面对的并不是一笔可以从北京直接拨付的总账：正额之外的附加征收散在各省州县，官员的日常供给也常倚赖这些不易核清的收入。西北的军情则不断提醒朝廷，银两在账上有了着落，并不等于粮食、牲畜和人手已经走到边地。雍正朝的整顿因此同时处理两个相连而不同的难题：谁能动用钱粮，谁又能使急令穿过层层衙门、道路和季节。

\begin{event}{\eventref{EQ-1726-HAOXIAN-GUIGONG}{耗羡归公与养廉银的安排}}{雍正四年（一七二六年）}{各省财政与州县官员}{制度方向可核；各地征收、留支和民间负担不能由一项安排概括}
耗羡归公所针对的，正是正额俸银不足而地方附加征收长期并存的局面。制度设计试图把原先分散的收入归入较可稽核的省级财政，再以养廉银供应官员；它希望把「为官须有用度」和「征收须可查问」放进同一套账目中。对督抚而言，这增加了汇总、报解和分配的责任；对知县、书吏和纳户而言，原先在地方关系中处理的名目、期限和折色，开始更多地进入省级与部院的核算。

不过，制度并不会自行消除中间环节。省里如何分拨，县里怎样征收，胥吏、乡绅、里役和纳户在何处协商或承受费用，都要在不同地方的账册、案件和文书中追问。归公可以说明朝廷试图改变财政的可见性，不能据此断定每一县的附加负担已经减少，也不能把养廉银写成地方官从此不再依赖其他资源。改革把责任链拉得更长，责任能否真的追到末端，仍取决于地方的人和账。
\end{event}

三年后，西北用兵的紧迫性又使文书路径成为制度问题。军机房起初处理的是紧急军政事务，后来发展为军机处；它缩短了皇帝、近臣与部院之间的呈递和批答线路，却没有把前线的距离一并缩短。命令送出之后，还要由将领、督抚、驿递人员、运粮者和地方承办者把它转成可行的行动。因而，军机处的意义不在于「中央从此无所不至」，而在于它让危机中的信息和决断更集中地流动；等待粮草、道路不通或地方难以筹集夫役时，这条流动仍会停在纸面与沿途之间。

\begin{event}{\eventref{EQ-1729-JUNJICHU-FOUNDING}{设军机房，处理紧急军政文书}}{雍正七年（一七二九年）}{北京与西北军政网络}{设立及文书处理方向可核；命令抵达、地方筹办和前线结果须分别检验}
军机房的出现不是抽象的「集权」标志，而是常规部院程序面对西北战争时的一次加速。更快的批答可以把调兵、筹饷和问责连在一起，也使北京较容易看见来自边地的呈报；但呈报本身已经经过地方官、将领和翻译链条的筛选。前线所需的是可到达的粮械、可通行的道路和愿意合作的人，而不是一份已经朱批的文书。

这套机构后来成为中后期决策的枢纽，却始终要和省级财政、驿路、军台以及地方中介共同行动。把军机档里连续的发令当作连续的执行，会抹掉那些没有及时到达、到达后难以办理，或在办理中改变了形状的环节。
\end{event}

乾隆初年的金川战事把这层差距摆在山路上。乾隆三年，清军对大金川展开大规模干预；山地、堡寨和土司关系并不允许远征按北京预定的节奏推进。到次年，粮运、道路和军费已成为持续的难题：粮食、牲畜和炮械必须经过有限的转运线，驻军越久，四川及邻近地区的筹办便越多地卷入战场。这里的成本不只是一笔军费，还包括道路的占用、民夫的征调，以及村社和市场为远征让出的劳力与物资。

乾隆十年，朝廷尝试以招抚、盟约和军事压力重组当地关系；次年以和议和行政安排暂告收束。它是推进困难和军费高昂条件下的阶段性选择，不是地方权力冲突已经消失的证明。第一次战争留下的土司关系没有被完全改写，二十余年后再度成为冲突条件。由此可见，远征的结局从来不止在攻守之间决定，也在谁能供粮、谁能通路、谁愿意接受新的关系之中决定。

同一时期，内地的财政与交通也没有因为边地用兵而暂停。黄河堤防、河道和经费的奏报与工程，关乎沿河社区，也关乎漕运和田赋能否维持；定额写在纸上，河势、工料和民工的实际处境却不能由定额推出。乾隆十四年，江南钱粮、漕运和士绅关系仍在户部与地方之间反复讨论，正说明富庶地区的税粮如何征解，始终是全国财政的一部分。盐政的整顿亦然：引岸、课额与商人网络相互牵连，朝廷能够调整规制，却不能仅凭一项整顿就断言盐价、私盐或商人的收益已经随之改变。

八旗的供养压力又把财政问题落到家庭的身份上。雍正五年编修谱牒，意在界定旗籍、世系和待遇资格；乾隆七年允许部分汉军旗人出旗，乾隆二十一年又命部分副户口改归民籍。每一步都显示国家试图重新排列供养资格与财政资源，但对当事家庭而言，身份变化还要经过登记、迁居、谋生和地方接纳。名册能够说明分类的意图，不能替代这些家庭如何获得生计的记录。国家分类在这里既是财政技术，也是会改变家庭选择空间的社会过程。

\subsection{扩张不是地图上的一次涂色}

西北的扩张同样不能从一张完成后的地图倒读。乾隆十九年，阿睦尔撒纳等准噶尔贵族转向清廷，改变了联盟格局；次年清军分路进入伊犁，击溃达瓦齐政权并占领准噶尔核心地区。但一七五五年后秩序并未固定：阿睦尔撒纳的反清使冲突再起，随后天山南路的战争又把北路与南路的行军、补给和地方关系连成新的难题。一七五八至一七五九年，清军进入喀什噶尔、叶尔羌等地，朝廷后来用「平定」概括这段过程；这个完成式只能说明朝廷如何命名战果，不能替代多条道路上反复发生的联盟变化、供给困难与暴力后果。

\begin{event}{\eventref{EQ-1760-ILI-GENERAL}{设伊犁将军，筹办驻防、屯田与安置}}{乾隆二十五年（一七六〇年）}{伊犁与北疆}{设官和军政筹办方向可核；迁入人数、屯田成效及各群体处境不能由官文书直接量化}
战后设置伊犁将军，是把临时远征转成长期军政管理的尝试。驻防、屯田、道路和市场必须一起运作：若粮食、种子和劳力不能持续到位，驻军便难以自给；若交通不能连接伊犁与东面的节点，调度仍要依赖漫长而脆弱的补给线。乾隆二十七年迪化等北疆城市的军政建设加快，成为屯田与交通的枢纽；次年继续安置屯田人口、军户和商业网络，正是要使守备不只停留在营垒之内。

这不是把一片空地填入帝国版图。旗兵及其家属、绿营、屯田者、商人以及原有的地方社会，都在不同的条件下面对土地、水利、粮饷和市场。乾隆二十九年部分八旗兵丁调往西北，边务也随之进入家庭的迁徙与供给。户籍和屯田册能显示国家希望谁在哪里定居，不能自动说明谁留下、怎样生活，或彼此如何相处；这些问题需要不同语言的地方文书、道路与军需材料留下的零散痕迹来约束。
\end{event}

第二次金川战争使西北与西南共同显出扩张的财政边界。乾隆三十三年再战开始后，山地堡寨使通常的野战难以奏效；次年，道路、炮兵、粮运和将领调度不断加码。乾隆三十五年，四川及邻省为军需转运粮食、牲畜并征调民夫，战场之外的道路和村社因而承受压力。朝廷最终在乾隆四十一年攻克金川，次年推行改土归流、设治和驻防，意在把军事胜利变成持续的行政安排。可是，土地、劳役和宗教网络如何重组，并不会随着设治的名称立即完成；乾隆四十三年地方官仍把四川钱粮、道路与民力恢复列为要务，正说明战后的恢复本身也是一段过程。

\subsection{禁令与账册抵达地方时}

\begin{event}{\eventref{EQ-1724-CATHOLIC-PROSCRIPTION}{扩大对天主教活动的限制}}{雍正二年（一七二四年）}{清廷、地方官与天主教社群}{限制政策可核；不同地区教民、传教士和官员的处境及执行程度差异很大}
一七二四年的限制把传教、地方教民和外来人员放入治理问题。礼仪争论与地方官对宗教网络的疑虑在中央相遇，形成了禁令的方向；它并没有规定每一座城、每一个村必须以同一种方式处理教民。地方官还要在户口、治安、宗教关系和对外往来之间作判断，教民、传教士、地方士绅和邻里也不是被动地等待一纸命令抵达。

到雍正十二年，地方官仍在执行限制，但不同地区的活动与冲突呈现不同的轨迹。官府文书可以让我们看见禁令、呈报和处分的语言，教会材料可以留下另一面的遭遇和网络；任一类材料都不能单独变成全国宗教生活的总图。较稳妥的读法，是把政策、地方执行、协商或冲突及其后果分开，而不把官府或教会的分类直接变成群体的本质。
\end{event}

\subsection{材料能看见什么，不能替谁作结论}

这些改革和远征留下的材料多从行动的上端开始：朱批奏折、户部奏销、军机档和方略最擅长保存呈报、批示、调度、奖惩以及朝廷希望人们如何理解战果。它们能够帮助读者分清一项办法何时提出、谁奉命筹办、北京何时收到消息，却不能直接替我们回答一个县多收了什么、一条山路实际走了多少人，或一户迁来的家庭是否安稳定居。

要追问这些问题，需要把中央文书同地方账簿、道路和军需记录、契约、碑刻、社区材料以及满文、蒙古文、维吾尔文、俄文等不同语境的记录并读。材料的缺口并不许我们用想象补满：金川的民夫、伊犁的迁徙、江南的纳户和各地教民，都不会在同样密度的档案中出现。也正因此，「改革成功」「疆域既定」或「禁令严行」都只能是待检验的概括。高峰时期的国家确实扩大了调度账册、命令和军队的能力；它同时把地方社会更深地卷入征收、转运、身份重整与战后恢复，而这些成本和回应，仍须在各地不完整的证据中谨慎辨认。

\section{伊犁之后：征服不是均质的治理}

1755年清军进入伊犁、击溃达瓦齐；阿睦尔撒纳旋即反清，1758至1759年南路大小和卓政权又告覆亡。战报可以把这些行动排成一条向西推进的胜利线，却不能使天山南北在同一时刻、以同一方式进入清廷的治理。伊犁河谷的驻防、军台、屯田与安置，面对的是牧地、水源和越冬条件；南疆绿洲的伯克、渠水管理者、城镇商人和农户，则在税役、市场与灌溉的安排中面对新政权。将军衙门可以调遣旗人、满汉军人与屯户，也必须依赖地方掌握道路、语言、收获时令和水利的人。征服改变了这些关系，却没有把它们抹平。

供给把这一区别变成日常的政治问题。军粮要沿军台和商路转运，牲畜需有草场，驻军与来迁者又同既有居民争取水、地和贸易机会；一封要求催粮或开垦的奏折，只能证明上级看见了短缺并试图处置，不能证明粮食准时抵达，或某一条渠水的分配已被接受。乌什在1765年的动荡因而不只是驻军体系里的一次“意外”：征发、劳役、地方权威与绿洲生活相互挤压，足以使一座有城、有兵的地方仍然失序。清廷随后能够镇压并重组行政安排，但镇压的完成与社会关系的稳定是两件不同的事。

西藏、青海、尼泊尔边境、川西山地和东南海疆也各有自己的瓶颈。驻藏大臣与噶厦关系的调整、廓尔喀战争后的规条，确实把更多事务接到北京的边务文书中；消息和物资却要穿过寺院网络、山口、驿道和地方政教关系。清缅战争、安南撤兵与台湾林爽文起事所显出的，并非一个简单的“边远地区难以控制”，而是病疫、雨季、航道、翻译、居民合作和外交礼仪会在不同地方以不同方式限制军令。帝国可在数个方向同时用兵，却不能让每一条道路、每一座港口和每一个中介都按中枢的时间表运转。

广州一口通商与马戛尔尼使团把同样的问题推到港口。行商、海关官员、水师、闽粤商人和抵达的欧洲船只并非站在一条命令链上；货物、欠款、引水、停泊与礼仪的规则须在具体交涉中落定。把一口通商写成闭关不动，或把使团写成中国已经“开放”的开端，都会遗漏这些在港口实际承担风险的人。至白莲教战争爆发时，高额军费、山区道路、地方结社与动员已重新缠结起来：扩张并未只留下疆域，也留下需要由户部、地方官、民夫、商人和村寨继续承担的财政与治理压力。

\subsection{材料与争议}

满汉《实录》、军机处档、理藩院档与各类方略最适合追索任命、调兵、赏罚和清廷如何叙述战事；它们让人看见一项命令进入了何种机构，却常把抵达后的协商、拖延与损失压缩为结果。地方奏报同样是官员为考成、请饷或避责而写的陈述。军报中的歼获、伤亡和“平定”因此应当首先读作战争中的国家报告，而不是可以直接累加为人口后果或地方态度的统计。

维吾尔／察合台文、蒙古文和藏文文书，寺院记录、契约与地方档案，可在一些地点补出水利、土地、捐施、贸易和地方机构的时间；俄国、尼泊尔、缅甸及欧洲来往文书又保存了跨境行动者的另一种观察位置。这些材料的留存并不平均。它们不能合成一幅无缝的“地方声音”，正因为如此，缺少自述的社群不应由官方分类或后出的民族名称替其发言。能较稳妥地说的是某项命令被下达、某次战事被报告，或某份契约记录了一项关系；至于服从的范围、迁徙的规模和暴力的总量，仍须按地点、语言和文类分别追问。

\subsection{伊犁、喀什噶尔与绿洲的水}

伊犁将军衙门的文书把驻防、屯田、军台和赏罚编入征服后的秩序。它们可确认清廷试图把兵、粮、土地和人事联结起来，却不能把伊犁河谷的牧地与南疆城镇的渠水当作同一种可调度资源。前者牵连牧群的季节移动与驻防需求，后者牵连绿洲耕作、城镇市场和地方权威；商队也要在关卡、治安和跨境机会之间判断行程。阿睦尔撒纳、大小和卓、伯克、驻防官和商人处于不同的联盟与风险之中，不能被概括为“归附者”与“反叛者”两类。

乌什的情形使这种差别格外清楚。军粮能够进入城中，只说明转运在某一刻成功；它并不回答由谁交纳、谁服役、谁裁断税役与灌溉的纠纷。把《平定准噶尔方略》的凯旋叙述同满文军机档、维吾尔／察合台文契约以及俄国和中亚材料并读，才可在局部重建水、地、贸易与迁徙怎样被重排。即使如此，保存下来的文书也更容易记录官府、商人与能订立契约者，较难保存佃户、妇女或临时迁徙者的选择；不能从档案的沉默推断他们没有承担代价。

\subsection{拉萨、金川与西南山道}

清廷在拉萨调整驻藏大臣及相关规条时，面对的不是一块可由一纸章程均质管理的高原。噶厦、寺院、青海蒙古诸部、廓尔喀使团和驿道上的运输者分别影响消息与供给能否越过山口。藏文寺院文书能够显示某些机构、捐施和地方时间，清廷档案则能够显示驻军、任命与奏议；二者都不能独自代表整个高原，更不能由北京的制度分期替代当地政治的节奏。

川西金川把财政成本压在山道的长度和坡度上。土司关系、碉楼之间的通道、民夫粮运与季节性疾病，使清军即使取得一处据点，也要继续为炮械、粮食和人员的抵达付出代价。所谓改土归流的结果，须由地方档、碑刻、诉讼与迁徙记录追问土地、债务和官署究竟在哪些村寨改变；战报中的“平定”不能预先回答这些问题。金川的战后安排也不能替高原或西南其他地区作结论，正如驻藏体制不能替山地道路解决供给。

\subsection{广州、台湾与海上通路}

广州的行商、海关官员、福建与广东水师、闽南商人与欧洲船只共同限定贸易和防务的边界。一口通商调整了合法贸易的门槛，却没有让海岸社会停止流动：船只的航期、港口的引水、商人的信用和沿岸居民的生计仍使规则在执行中不断被重谈。海关账册和航海日志可以说明一条航线、一批货物或一次交易的制度位置，不能由此推出整个海岸都遵从同一秩序。

1787年林爽文起事的军报记录清军收复城镇，台湾府县志、契约与平码材料则把垦殖、租佃、渡海和地方武装放回具体地点。平原村落、鹿港等港口与内山社群承受的征发和接触，不能被一个“收复”或“设治”概括。原住民社群的自述保存尤其稀少；海防命令、移民契约和殖民或官府记录都带着各自的土地与治理诉求。这里的证据缺口不是可以用更响亮的国家叙事填补的空白。

\section{读者事件簇：乌什的驻防与金川的山道}

\eventref{EQ-1765-USH-REBELLION}{乌什事件（1765）}发生在征服后的南疆驻防体系之内。乾隆三十年的实录、军机处档和战事方略能跟出清廷接到消息、调兵和处理善后的次序；它们的用语也首先服务于镇压与重组治理的需要。若只沿着这条文书链阅读，城内税役、劳役与地方权威的冲突就会被化为抽象的“不服”。维吾尔和察合台文材料、地方文书与跨境商路记录虽各自残缺，却提醒读者，军粮入城和官员到任都不足以证明当地社会已经接受新的权力关系。

\eventref{EQ-1776-JINCHUAN-CONQUEST}{第二次金川战争结束（1776）}把类似的限度放在川西山道上。攻克碉楼结束了一轮军事行动，却没有终止人力和物资的账目：夫役、道路、炮兵以及四川和邻省的转运，已把负担分散给许多并不出现在凯旋叙述中的村寨和家庭。战后设治、土地关系与迁徙的变化需要另以地方材料检验。乌什与金川并非同一场冲突；前者使绿洲驻防的摩擦可见，后者让山地后勤的延续显形。两者共同阻止读者把驻防城或“平定”误读为已经建立了同一种秩序。

\section{读者事件簇：青海部署与高原中介}

\eventref{EQ-1728-ZHUNGHAR-TIBET-SETTLEMENT}{调整驻藏与青海蒙古部署（1728）}，是在拉萨政局和草原力量关系不稳、准噶尔威胁仍在的条件下作出的军政回应。雍正六年的实录、理藩院和相关军政文书可以定位清廷的部署，却不能代替高原上的实施经过。驻防名额写在文书中以后，消息仍要经由青海蒙古诸部、寺院网络、驿道与运输者传递；粮草能否前行，也取决于山口、季节和地方协商。清廷由此更直接地协调西藏事务，并不意味着它预先决定了这些中介会在何处配合、拖延或改写命令。

\section{读者事件簇：乾隆继承的不是空白账册}

\eventref{EQ-1735-QIANLONG-ACCESSION}{乾隆即位（1735）}时，雍正朝的财政整顿、朱批奏折与军机文书已使许多事务能够较快地进入中枢。雍正十三年的实录、奏折与题本能使人看见继承完成和官僚程序的延续，却不能把档案的密集误作地方问题已经被掌握或解决。新皇帝继承的是可以调度的制度，也继承了等待核销的亏空、需要转运的粮饷、山地与边地的供给难题，以及必须由地方官、商人和社群承担的执行成本。此后的扩张正是在这些并未清空的账目上展开。

\section{读者事件簇：第一次金川战争的山道}

\eventref{EQ-1738-FIRST-JINCHUAN-WAR}{大金川战事展开（1738）}时，清廷对川西治理的扩展卷入既有的地方权力冲突。乾隆三年的战报与地方奏报保存了进军、请饷和处置的官方序列；碉楼之间狭窄而曲折的通道、土司间的关系和四川粮运，却使兵力进入山谷不等于控制道路。清军后来以招抚与和议收束战事，显示军事推进必须同地方联盟和供给的限度谈判。不能从一份讲述行动的奏报推出山地社会已经按清廷设想重新排列。

\section{读者事件簇：第二次金川战争的夫役账}

\eventref{EQ-1768-SECOND-JINCHUAN-WAR}{第二次金川战争开始（1768）}，意味着第一次和议未能消除土司与清廷之间的权力冲突。炮兵、粮草和民夫须从四川及邻省沿山道转运；每一次征发或加派都可能暂时接上前线，却把劳力、债务与村寨缺人的风险留在后方。军机处档和战事方略能追踪更大规模资源投入的意图与报告，难以单独计量劳役如何落在不同家庭。战争的代价因而不能只由攻下几座碉楼来计算，它也为1776年之后的设治与迁徙留下了需要继续处理的条件。

\section{读者事件簇：土尔扈特回归的安置难题}

\eventref{EQ-1771-TORGHUT-RETURN}{土尔扈特部回归（1771）}把伏尔加草原、俄国边境与伊犁的牧地安排接到清廷面前。实录、军机处和理藩院档案记录接纳、赏赐与安置的军政语言，也显示这不是单纯的边界新闻：部众迁徙之后，水草、牲畜、疾病、驻防与邻近部落的资源边界都要重新协商。俄国文献、蒙古文材料与清廷档案对迁徙的叙述位置不同，不能由任何一方的数字概括全部损失或内部差异。“安置”不是行政程序的终点，而是草场、军政和跨境关系中一段新的不确定开端。

